\begingroup
\section{The signed return through the actual adjacent relation and quotient}
\label{sec:sxr}

Retain the full original hypothetical packet
$\rho=1/2+\delta+i\gamma$, $0<\delta<1/2$, $\gamma>2$, and its
complete order $m$. Set
\begin{equation}\tag{SXR1}
\begin{gathered}
 g=2\xi,\quad h(s)=\prod_{z\in\{\rho,\bar\rho,1-\rho,1-\bar\rho\}}
                    (s-z)^m,\quad v_h=g/h,\\
 w_h(y)=|v_h(1/2+iy)|^2/(2\pi),\quad m_k=w_h^{*k},\quad
 k=4l+1\ge5,\quad e=1+k(m-1),\\
 q=e(k+1)^2,\quad c=k/2,\quad
 \chi(S)=\prod_{a,b=0}^k
  (S-c-(2a-k)\delta-i(2b-k)\gamma)^e.
\end{gathered}
\end{equation}
The full Taylor unit, source masses, period parameter and tensor
coordinates remain fixed. All limits fix this complete packet and
use the displayed admitted tensor orders. The original zero hypothesis
is the one in the total counterfactual; no numerical offcritical zero
is supplied or asserted by this calculation.

\subsection{The unchanged source path and its exact derivatives}

The original CA14a source interpolation is, with its parameter written
$x$ to distinguish it from the period $u$,
\begin{equation}\tag{SXR2}
 d\mu_x(y)=(1-x)d\sigma(y)+x m_k(y)dy,\quad
 d\sigma(y)=\frac{|\Gamma(1/4+iy/2)|^2}{2\pi}dy,
 \quad 0\le x\le1.
\end{equation}
In particular its mass is
$(1-x)\sqrt{2\pi}+x(\int w_h)^k$. There is no replacement by
a probability measure. On $\mathcal P_d=\mathbb C[S]_{\le d}$
use the basis $1,S,\ldots,S^d$ and
\[
 (H_d(x))_{ab}=\int\overline{(c+iy)^a}(c+iy)^b\,d\mu_x(y),
 \qquad 0\le a,b\le d.
\]
These positive matrices have $d+1$ rows. Embeddings into
$\mathcal P_{2q}$ are literal zero extensions. Write
\begin{equation}\tag{SXR3}
 M(x)=H_{2q}(x),\quad \dot M=H_{2q}^{\rm ar}-H_{2q}^{\sigma},
 \quad C(x)=M(x)^{-1}\dot M.
\end{equation}
The dot denotes this source derivative only. The operator $C$ is
self-adjoint in $M$, since $MC=\dot M=\dot M^*$. For any two
polynomials in the ambient space,
\[
 \langle P,CQ\rangle_x
   =\int\overline{P(c+iy)}Q(c+iy)
                     (m_k(y)-\sigma(y))dy.
\]
This identity proves its exact connection to the original density
change. It is not a freely chosen coefficient operator.

Let $p_d(x)$ be the original monic source orthogonal polynomial and
$u_r(x)$ the monic polynomial in the original relation norm. Thus
\begin{equation}\tag{SXR4}
 \omega_d(x)=\|p_d(x)\|_x^2,\quad
 \nu_r(x)=\|\chi u_r(x)\|_x^2,\quad
 T_r(x)=\frac{\nu_r(x)}{\omega_{q+r}(x)},\quad d=q+r.
\end{equation}
The polynomial $u_r$ is orthogonal to degrees below $r$ for the
literal measure $|\chi(c+iy)|^2d\mu_x(y)$. Its full $S$-coordinate
monicity is retained. Positive finite moment Grams give unique monic
minimizers, and solving their linear systems shows that their
coefficients are smooth rational functions of $x$. Denominators
are positive on the whole compact interval. Differentiating the
minimum norms gives exactly
\begin{equation}\tag{SXR5}
 \dot\omega_d=\langle p_d,Cp_d\rangle_x,\qquad
 \dot\nu_r=\langle\chi u_r,C\chi u_r\rangle_x.
\end{equation}
Indeed $\dot p_d$ has degree below $d$ and $\chi\dot u_r$ lies
in $\chi\mathcal P_{r-1}$. Both coefficient-derivative cross terms
vanish by their respective defining orthogonalities. This proves
SXR5 without differentiating an asymptotic remainder.

The source is even. In real $y$ coordinates its monic orthogonal
polynomial has simple real roots by the sign-change proof in TCN4.
The polynomial $\chi(c+iy)$ has nonreal roots, so $\chi u_r$ cannot
be that monic minimizer. Therefore $T_r(x)>1$ for every $x$ and
$0\le r\le q$. The phase change is exactly
$P(S)\mapsto i^{-d}P(c+iy)$ in degree $d$; no source norm changes.

\subsection{An orthogonal exchange in the actual quotient section}

Put $b_d(x)=[p_d(x)]_\chi$ in the original quotient
$E=\mathbb C[S]/(\chi)$ and keep
\[
 G_d=(J_dH_d^{-1}J_d^*)^{-1},\qquad R_d=H_d^{-1}J_d^*G_d.
\]
For $d=q+r$ define the following actual source vectors:
\begin{equation}\tag{SXR6}
\begin{gathered}
 f_r=\frac{p_d}{\sqrt{\omega_d}},\quad
 g_r=\frac{\chi u_r}{\sqrt{\nu_r}},\quad
 \tau_r=T_r^{-1},\quad
 z_r=\frac{b_d}{\sqrt{\omega_d(1-\tau_r)}},\\
 w_r=\frac{f_r-\sqrt{\tau_r}g_r}{\sqrt{1-\tau_r}}
           =R_d z_r.
\end{gathered}
\end{equation}
All square roots here are the positive square roots of the original
positive real norms. To prove the equalities, monicity and the
degree-$d$ orthogonality give
$\langle p_d,\chi u_r\rangle_x=\omega_d$, hence
$\langle g_r,f_r\rangle_x=\sqrt{\tau_r}$ with this exact positive
phase. The complete relation space in degree $d$ is
$\chi\mathcal P_r$. Both $p_d$ and $\chi u_r$ are orthogonal to
$\chi\mathcal P_{r-1}$, and the coefficient of the last orthogonal
relation vector in $p_d$ is $\omega_d/\nu_r$. Consequently
\begin{equation}\tag{SXR7}
 R_db_d=p_d-\frac{\omega_d}{\nu_r}\chi u_r,
 \qquad b_d^*G_db_d=\omega_d(1-T_r^{-1}).
\end{equation}
This is the actual full relation projection, not a deletion of its
highest row. It proves SXR6 and
$\|g_r\|_x=\|w_r\|_x=1$, $\langle g_r,w_r\rangle_x=0$.
The vector $g_r$ belongs to the original relation space and $w_r$
to the original minimum quotient section in the same degree.

For a unit vector $v$ write $P_v=vv^*M$ for its source-orthogonal
projector. SXR5 gives the exact identity
\begin{equation}\tag{SXR8}
 \frac{d}{dx}\log T_r
   =\operatorname{Tr}((P_{g_r}-P_{f_r})C)
   =\langle g_r,Cg_r\rangle_x-\langle w_r,Cw_r\rangle_x
      +\operatorname{Tr}((P_{w_r}-P_{f_r})C).
\end{equation}
In the orthonormal basis $(g_r,w_r)$, the last projector difference
is exactly
\[
 P_{w_r}-P_{f_r}=
 \begin{pmatrix}
 -\tau_r&-\sqrt{\tau_r(1-\tau_r)}\\
 -\sqrt{\tau_r(1-\tau_r)}&\tau_r
 \end{pmatrix}.
\]
Its eigenvalues are $\sqrt{\tau_r},-\sqrt{\tau_r}$; all remaining
ambient eigenvalues are zero. If
$\operatorname{osc}_x C=\lambda_{\max}(C)-\lambda_{\min}(C)$,
the Rayleigh quotients at the two unit eigenvectors therefore prove
\begin{equation}\tag{SXR9}
 |\operatorname{Tr}((P_{w_r}-P_{f_r})C)|
          \le T_r^{-1/2}\operatorname{osc}_x C.
\end{equation}
Every scalar part of $C$ cancels exactly because this projector
difference has trace zero. A bound by an uncentered source norm is
neither needed nor substituted for the oscillation.

\subsection{All four endpoint signs before the estimate}

Let $V_N(x)=\det G_N(x)$ and retain precisely
\begin{equation}\tag{SXR10}
\begin{gathered}
 \mathcal B_k(x)=\log V_{q-1}(x)+\log V_q(x)
                     -\log V_{2q-1}(x)-\log V_{2q}(x),\\
 \delta_k^\sigma=\mathcal B_k(1)-\mathcal B_k(0),\qquad
 c_0=c_q=1,\quad c_r=2\ (1\le r\le q-1).
\end{gathered}
\end{equation}
The monic source/relation change of basis has determinant one, so
$V_{q+r-1}=\prod_{d=0}^{q+r-1}\omega_d/\prod_{j=0}^{r-1}\nu_j$.
Taking consecutive ratios and then the four displayed endpoint signs
gives the exact finite equality
$\mathcal B_k(x)=\sum_{r=0}^q c_r\log T_r(x)$.
In particular $\sum c_r=2q$, with coefficient one at both endpoints.

Define the original relation-versus-quotient receiving integral
\begin{equation}\tag{SXR11}
 \mathcal Q_k=\int_0^1\sum_{r=0}^q c_r
  \bigl(\langle g_r,Cg_r\rangle_x
                  -\langle w_r,Cw_r\rangle_x\bigr)dx.
\end{equation}
This uses the original source path, actual minimum sections, full
relation rows and positive monic phases of SXR6. It is an explicitly
constructed scalar of that source. SXR8 gives
\begin{equation}\tag{SXR12}
 \delta_k^\sigma=\mathcal Q_k+\mathcal E_k,\qquad
 \mathcal E_k=\int_0^1\sum_{r=0}^q c_r
          \operatorname{Tr}((P_{w_r}-P_{f_r})C)dx.
\end{equation}
There is no assertion that either summand has a necessary sign.

Let $b_-,b_+$ be the extreme eigenvalues of the unchanged relative
Gram $(H_{2q}^\sigma)^{-1}H_{2q}^{\rm ar}$. Diagonalizing its
self-adjoint representative in the original Gamma metric gives
the exact eigenvalues of $C(x)$:
$(b_j-1)/(1-x+xb_j)$. This expression increases in $b_j$,
with derivative $(1-x+xb_j)^{-2}>0$. Direct integration gives
\begin{equation}\tag{SXR13}
 \int_0^1\operatorname{osc}_x C\,dx
       =\log(b_+/b_-).
\end{equation}
The formula at $b_j=1$ follows by continuity. It preserves even a
large common source mass until the ratio cancels it exactly.

\subsection{A finite certificate and an exponentially small error}

Keep the original TW constants, without changing their values:
\[
 \frac{a_k}{D_d}\|P\|_\sigma^2\le\|P\|_{\rm ar}^2
                      \le A_k\|P\|_\sigma^2,
 \quad \deg P\le d,\quad
 L_{k,d}=\log(A_kD_d/a_k).
\]
Their full definitions are TCN10 and the unchanged TW/BT providers
in the source closure; $L_{k,2q}=O_h(k+\log q)$. On SXR2 they
give lower and upper factors
$l_d(x)=1-x+x a_k/D_d$, $u(x)=1-x+xA_k$.
The ratio $u(x)/l_d(x)$ increases from $1$ to
$A_kD_d/a_k$: its derivative has numerator $A_k-a_k/D_d\ge0$.
The last inequality follows already by applying TW to a nonzero
polynomial. Taking the actual monic minima in numerator and
denominator therefore proves, simultaneously for every $x$,
\begin{equation}\tag{SXR14}
 T_r(x)\ge e^{-L_{k,2q}}T_r^\sigma,
 \qquad 0\le r\le q,\qquad
 \log(b_+/b_-)\le L_{k,2q}.
\end{equation}
Here $T_r^\sigma$ uses the same complete $\chi$ and its monic
relation norm in the original Gamma source, not a replacement
arithmetic packet.

It is a finite explicit determinant quantity. With
$L_r^\sigma=\operatorname{Gram}_\sigma(\chi S^j)_{0\le j<r}$,
$\det L_0^\sigma=1$, define
\begin{equation}\tag{SXR15}
 \Theta_k=e^{-L_{k,2q}}
   \min_{0\le r\le q}
   \frac{\det L_{r+1}^\sigma}{\det L_r^\sigma}
       \frac{4^{q+r}}{\sqrt{2\pi}(2q+2r)!}.
\end{equation}
Every determinant here is in its original monic frame and retains
the full Gamma mass. For example its moments are obtained exactly
by differentiating
$\int e^{ity}d\sigma(y)=\sqrt{2\pi}(\cosh t)^{-1/2}$ at zero.
This is the same Gamma Fourier identity used in the retained source
provider, with its stated mass. No arithmetic zero is needed to
compute these reference matrices at fixed packet parameters.

Combining SXR9, SXR12--15 and the strict inequality $T_r(x)>1$
proves the finite bound
\begin{equation}\tag{SXR16}
 \boxed{|\delta_k^\sigma-\mathcal Q_k|
 \le 2q\min\{1,\Theta_k^{-1/2}\}\log(b_+/b_-)
 \le 2q L_{k,2q}\min\{1,\Theta_k^{-1/2}\}.}
\end{equation}
All $q+1$ adjacent exchanges and both original endpoint weights
were included before this estimate was taken.

For completeness the new TCN/TCX estimates evaluate the size of
this finite bound. They give
\begin{equation}\tag{SXR17}
 \min_{0\le r\le q}\log T_r^\sigma
        =2q a_1+O_h(k+\log q),\qquad a_1=\psi(1)>0.
\end{equation}
Here is the required uniform argument rather than an inference
from a four-endpoint assertion. For
$0\le r\le\lfloor q/64\rfloor$, TCX2--9's original-root interval
and monic arcsine bound apply to $\sigma$ with comparison factor
one. They give a lower rate $c_*q$, $c_*>1/3>2a_1$, with
error $O_h(k+\log q)$. For the remaining integers, TCN13 and
TCN17--26 give $\log T_r^\sigma=2q\psi(r/q)+O_h(k+\log q)$
uniformly on the fixed compact ratio interval $[1/128,1]$,
retaining the original packet cutoff. Since $\psi$ decreases,
this is at least $2qa_1-O_h(k+\log q)$. At $r=q$ the same formula
gives a matching upper bound for the minimum. The original cutoff
holds eventually because $q\ge(k+1)^2$ and
$R=k\sqrt{\delta^2+\gamma^2}$. This proves SXR17 and
$\log\Theta_k=2qa_1+O_h(k+\log q)$.
Consequently the exact error in SXR12 satisfies
\begin{equation}\tag{SXR18}
 \boxed{|\mathcal E_k|
       \le\exp\{-qa_1+O_h(k+\log q)\}=o(q).}
\end{equation}
This is a bound on an actual signed correction term, not a claimed
evaluation of the still-present $\mathcal Q_k$.

\subsection{Return to the original joint kernel and observable image}

At fixed original nonzero period and branch, keep the complete
observation $O:E\to\mathbb C^b$ onto its fixed image frame, with
its full Taylor units, reduced mixed map and ordered periods. It
does not vary with the source interpolation. Let $I_K$ be the fixed
kernel frame. At each actual degree $d=q+r$ set
\begin{equation}\tag{SXR19}
\begin{gathered}
 Q_d=(OG_d^{-1}O^*)^{-1},\quad
 P_{K,d}=I_K(I_K^*G_dI_K)^{-1}I_K^*G_d,\\
 P_{Y,d}=G_d^{-1}O^*Q_dO,\quad
 w_{K,r}=R_dP_{K,d}z_r,\quad w_{Y,r}=R_dP_{Y,d}z_r.
\end{gathered}
\end{equation}
Empty kernels or images use their unique empty matrices. These
are the original source-Gram minimum sections, at their written
degrees; no boundary metric is substituted. Multiplying the
formulas gives $OP_{K,d}=0$, $OP_{Y,d}=O$ and orthogonality of
their ranges in $G_d$. Hence $P_K+P_Y=I$ and
$w_r=w_{K,r}+w_{Y,r}$, with the two summands orthogonal in $M(x)$.
The literal receiving integrand is therefore
\begin{equation}\tag{SXR20}
\begin{aligned}
 \langle g_r,Cg_r\rangle_x-\langle w_r,Cw_r\rangle_x
 ={}&\langle g_r,Cg_r\rangle_x
   -\langle w_{K,r},Cw_{K,r}\rangle_x
   -\langle w_{Y,r},Cw_{Y,r}\rangle_x\\
 &-2\Re\langle w_{K,r},Cw_{Y,r}\rangle_x.
\end{aligned}
\end{equation}
Orthogonality for $M$ does not remove the last pairing with $C$.
It is retained, with its original sign, inside $\mathcal Q_k$.
The quotient coefficient is exactly $z_r$ in SXR6, and its observed
value is exactly $Oz_r$. There is no assertion that it is the
distinguished primary or that its observable fraction is uniformly
nonzero. On the primary, the earlier nonvanishing subsequence and
the undivided JQC identity retain their proved domains.

The relation $\chi u_r$ has its original ordered theta primitive
at $W_*$ or a containing source. Its proper-source class and full
$h^2$-jet keep the TCD3--5 map and image restrictions. Applying the
linear source maps coefficientwise retains every coefficient face,
outer source label and nonempty leg mask. A zero remains at its
present support and is not the external $\tau$. Thus the controlled
exchange has been returned to the original mixed construction with
its exact morphisms; it has not erased that construction's kernel.

Equations SXR12, SXR16 and SXR20 leave a precisely specified original
mixed scalar with an exponentially bounded exchange error. The sign
and order-$q$ value of this scalar have not been computed here.
Neither the new error estimate nor the separate positive exterior
energy proves RH or certifies an offcritical zero.
\endgroup


\begingroup
\section{The surviving original mixed connection matrix}
\label{sec:CMX}

The matrix constructed here is the cross matrix of two specified
polynomial subspaces in the original source interpolation. Its entries
retain all monic norms and every arithmetic coefficient. Its diagonal
becomes exponentially small on the original degree window, while its
